\section{Introduction}
\label{sec:introduction}

Coastal lagoons concentrate fisheries and biodiversity along roughly 13\% of
the global shoreline \citep{Carrasco2016, Kjerfve1994}, and their shallow,
semi-enclosed geometry makes water exchange and renewal sensitive to
atmospheric and tidal forcing. In the Mediterranean, where tidal amplitudes
rarely exceed 0.3~m, wind is a significant driver of horizontal exchange, and
through it of the residence times and salinity gradients that condition benthic
habitat \citep{Niedda2007, Umgiesser2014}. In many of these lagoons that
habitat is seagrass, principally \emph{Posidonia oceanica} (L.) Delile and
\emph{Cymodocea nodosa} (Ucria) Ascherson, whose canopies attenuate incident
wave energy \citep{Infantes2012} and extract momentum from the mean flow
through a drag that modifies near-bed turbulence and the vertical velocity
structure \citep{Nepf2012}. Global seagrass extent has fallen by 19.1\% of the
surveyed area since 1880 \citep{Waycott2009, Dunic2021}, with the fastest
recent losses where meadows are exposed to coastal development and degraded
water quality \citep{Turschwell2021}. How circulation sets dispersal and
connectivity in vegetated lagoons therefore concerns transport pathways, a
quantity that water level, the usual validation target of lagoon models, does
not measure directly.

The \emph{Stagnone di Marsala}, on the north-western coast of Sicily, is a
lagoon of this kind, with a mean depth of about 0.95~m and a floor largely
covered by meadows of \emph{P.\ oceanica} and \emph{C.\ nodosa}
(Section~\ref{sec:study_area}), whose atoll and cordon formations lost 75\% of
their extent between 2003 and 2024 \citep{Maltese2025}.
\citet{LaLoggia2004} and \citet{DiMarca2009} characterised its tidal flushing,
residence times and turbidity with depth-averaged models, and
\citet{Ciraolo2009a} showed through particle tracking how strongly simulated
trajectories depend on the velocity field they are advected through.
\citet{DeMarchis2012} resolved the vertical and showed that wind drives a
recirculation with a near-bed return flow even at sub-metre depths, with the
meadow represented through Nikuradse roughness heights assigned by seagrass
class, and \citet{Ingrassia2024} calibrated a depth-averaged model against
water level for a range of vegetation scenarios, a setting in which an
equivalent bed roughness is the appropriate representation of the canopy.

Building on these studies, two aspects of the problem remain open for this
basin. The first is a canopy represented as a drag resolved over the vertical
extent of the plants, which acts on the shape of the velocity profile as well
as on its magnitude \citep{Nepf2012} and is now available in three-dimensional
solvers, together with the waves that alter bed stress in comparably shallow
Mediterranean lagoons \citep{Carniello2011, Morote2025}. The second is the
assessment of transport itself, the models of the basin having been evaluated
mainly on water level and on currents at fixed stations, whereas Lagrangian
observations are known to constrain a model differently from water level
\citep{Liu2011, Revelard2021}.

A three-dimensional model of the Stagnone is used here to separate the
contributions of the canopy and of wave coupling, with every combination of the
two scored against water level and observed drifter trajectories
(Section~\ref{sec:data_methods}). The study has two objectives, to rank the two
processes by their effect on Lagrangian transport skill, and to establish
whether water level, the quantity against which shallow lagoon models are
conventionally validated, identifies that ranking.
